\begin{register}{H}{LP\_GPIO\_OUT\_REG}{0x0004}\label{regdesc:LPGPIOOUTREG}
\regfield{(reserved)}{25}{7}{0000000000000000000000000}%
\regfield{LP\_GPIO\_OUT\_DATA\_ORIG}{7}{0}{{0x0}}%
\reglabel{Reset}\regnewline%
\begin{regdesc}\begin{reglist}
\label{fielddesc:LPGPIOOUTDATAORIG}\item [LP\_GPIO\_OUT\_DATA\_ORIG] Configures the output of GPIO0$\sim$GPIO6.\\
The value of each bit can be:\\
0: Low level\\
1: High level\\
Bit[0]$\sim$bit[6] are corresponding to GPIO0$\sim$GPIO6.\\ (R/W/WTC)
\end{reglist}\end{regdesc}
\end{register}


\begin{register}{H}{LP\_GPIO\_OUT\_W1TS\_REG}{0x0008}\label{regdesc:LPGPIOOUTW1TSREG}
\regfield{(reserved)}{25}{7}{0000000000000000000000000}%
\regfield{LP\_GPIO\_OUT\_W1TS}{7}{0}{{0x0}}%
\reglabel{Reset}\regnewline%
\begin{regdesc}\begin{reglist}
\label{fielddesc:LPGPIOOUTW1TS}\item [LP\_GPIO\_OUT\_W1TS] Configures whether or not to enable the output register \hyperref[regdesc:LPGPIOOUTREG]{LP\_GPIO\_OUT\_REG} of GPIO0$\sim$GPIO6.\\
The value of each bit can be:\\
0: Not set\\
1: The corresponding bit in \hyperref[regdesc:LPGPIOOUTREG]{LP\_GPIO\_OUT\_REG} will be set to 1\\
Bit[0]$\sim$bit[6] are corresponding to GPIO0$\sim$GPIO6.\\
Recommended operation: use this register to set \hyperref[regdesc:LPGPIOOUTREG]{LP\_GPIO\_OUT\_REG}. \\
(WT)
\end{reglist}\end{regdesc}
\end{register}


\begin{register}{H}{LP\_GPIO\_OUT\_W1TC\_REG}{0x000C}\label{regdesc:LPGPIOOUTW1TCREG}
\regfield{(reserved)}{25}{7}{0000000000000000000000000}%
\regfield{LP\_GPIO\_OUT\_W1TC}{7}{0}{{0x0}}%
\reglabel{Reset}\regnewline%
\begin{regdesc}\begin{reglist}
\label{fielddesc:LPGPIOOUTW1TC}\item [LP\_GPIO\_OUT\_W1TC] Configures whether or not to clear the output register \hyperref[regdesc:LPGPIOOUTREG]{LP\_GPIO\_OUT\_REG} of GPIO0$\sim$GPIO6.\\
The value of each bit can be:\\
0: Not clear\\
1: The corresponding bit in \hyperref[regdesc:LPGPIOOUTREG]{LP\_GPIO\_OUT\_REG} will be cleared.\\
Bit[0]$\sim$bit[6] are corresponding to GPIO0$\sim$GPIO6.\\
Recommended operation: use this register to clear \hyperref[regdesc:LPGPIOOUTREG]{LP\_GPIO\_OUT\_REG}.\\
(WT)
\end{reglist}\end{regdesc}
\end{register}


\begin{register}{H}{LP\_GPIO\_ENABLE\_REG}{0x0010}\label{regdesc:LPGPIOENABLEREG}
\regfield{(reserved)}{25}{7}{0000000000000000000000000}%
\regfield{LP\_GPIO\_ENABLE\_DATA}{7}{0}{{0x0}}%
\reglabel{Reset}\regnewline%
\begin{regdesc}\begin{reglist}
\label{fielddesc:LPGPIOENABLEDATA}\item [LP\_GPIO\_ENABLE\_DATA] Configures whether or not to enable the output of GPIO0$\sim$GPIO6.\\
The value of each bit can be:\\
0: Not enable\\
1: Enable\\
Bit[0]$\sim$bit[6] are corresponding to GPIO0$\sim$GPIO6.\\
 (R/W/WTC)
\end{reglist}\end{regdesc}
\end{register}


\begin{register}{H}{LP\_GPIO\_ENABLE\_W1TS\_REG}{0x0014}\label{regdesc:LPGPIOENABLEW1TSREG}
\regfield{(reserved)}{25}{7}{0000000000000000000000000}%
\regfield{LP\_GPIO\_ENABLE\_W1TS}{7}{0}{{0x0}}%
\reglabel{Reset}\regnewline%
\begin{regdesc}\begin{reglist}
\label{fielddesc:LPGPIOENABLEW1TS}\item [LP\_GPIO\_ENABLE\_W1TS] Configures whether or not to set the output enable register \hyperref[regdesc:LPGPIOENABLEREG]{LP\_GPIO\_ENABLE\_REG} of GPIO0$\sim$GPIO6.\\
The value of each bit can be:\\
0: Not set\\
1: The corresponding bit in \hyperref[regdesc:LPGPIOENABLEREG]{LP\_GPIO\_ENABLE\_REG} will be set to 1\\
Bit[0]$\sim$bit[6] are corresponding to GPIO0$\sim$GPIO6.\\
Recommended operation: use this register to set \hyperref[regdesc:LPGPIOENABLEREG]{LP\_GPIO\_ENABLE\_REG}.\\
(WT)
\end{reglist}\end{regdesc}
\end{register}


\begin{register}{H}{LP\_GPIO\_ENABLE\_W1TC\_REG}{0x0018}\label{regdesc:LPGPIOENABLEW1TCREG}
\regfield{(reserved)}{25}{7}{0000000000000000000000000}%
\regfield{LP\_GPIO\_ENABLE\_W1TC}{7}{0}{{0x0}}%
\reglabel{Reset}\regnewline%
\begin{regdesc}\begin{reglist}
\label{fielddesc:LPGPIOENABLEW1TC}\item [LP\_GPIO\_ENABLE\_W1TC] Configures whether or not to clear the output enable register \hyperref[regdesc:LPGPIOENABLEREG]{LP\_GPIO\_ENABLE\_REG} of GPIO0$\sim$GPIO6.\\
The value of each bit can be:\\
0: Not clear\\
1: The corresponding bit in \hyperref[regdesc:LPGPIOENABLEREG]{LP\_GPIO\_ENABLE\_REG} will be cleared\\
Bit[0]$\sim$bit[6] are corresponding to GPIO0$\sim$GPIO6.\\
Recommended operation: use this register to clear \hyperref[regdesc:LPGPIOENABLEREG]{LP\_GPIO\_ENABLE\_REG}.\\
(WT)
\end{reglist}\end{regdesc}
\end{register}


\begin{register}{H}{LP\_GPIO\_IN\_REG}{0x001C}\label{regdesc:LPGPIOINREG}
\regfield{(reserved)}{25}{7}{0000000000000000000000000}%
\regfield{LP\_GPIO\_IN\_DATA\_NEXT}{7}{0}{{0x0}}%
\reglabel{Reset}\regnewline%
\begin{regdesc}\begin{reglist}
\label{fielddesc:LPGPIOINDATANEXT}\item [LP\_GPIO\_IN\_DATA\_NEXT] Represents the input value of GPIO0$\sim$GPIO6.\\
Each bit represents a pin input value:\\
0: Low level input\\
1: High level input\\
Bit[0]$\sim$bit[6] are corresponding to GPIO0$\sim$GPIO6.\\ (RO)
\end{reglist}\end{regdesc}
\end{register}


\begin{register}{H}{LP\_GPIO\_STATUS\_REG}{0x0020}\label{regdesc:LPGPIOSTATUSREG}
\regfield{(reserved)}{25}{7}{0000000000000000000000000}%
\regfield{LP\_GPIO\_STATUS\_INTERRUPT}{7}{0}{{0x0}}%
\reglabel{Reset}\regnewline%
\begin{regdesc}\begin{reglist}
\label{fielddesc:LPGPIOSTATUSINTERRUPT}\item [LP\_GPIO\_STATUS\_INTERRUPT] Configures the interrupt status of GPIO0$\sim$GPIO6.\\
The value of each bit can be:\\
0: No interrupt\\
1: Interrupt is triggered\\
Bit[0]$\sim$bit[6] are corresponding to GPIO0$\sim$GPIO6.\\
This field is used together with \hyperref[fielddesc:LPGPIOPINNINTTYPE]{LP\_GPIO\_PIN\regindex{n}\_INT\_TYPE} in register \hyperref[regdesc:LPGPIOPINNREG]{LP\_GPIO\_PIN\regindex{n}\_REG}.\\
 (R/W/WTC)
\end{reglist}\end{regdesc}
\end{register}


\begin{register}{H}{LP\_GPIO\_STATUS\_W1TS\_REG}{0x0024}\label{regdesc:LPGPIOSTATUSW1TSREG}
\regfield{(reserved)}{25}{7}{0000000000000000000000000}%
\regfield{LP\_GPIO\_STATUS\_W1TS}{7}{0}{{0x0}}%
\reglabel{Reset}\regnewline%
\begin{regdesc}\begin{reglist}
\label{fielddesc:LPGPIOSTATUSW1TS}\item [LP\_GPIO\_STATUS\_W1TS] Configures whether or not to set the interrupt status register \hyperref[regdesc:LPGPIOSTATUSREG]{LP\_GPIO\_STATUS\_INT} of GPIO0$\sim$GPIO6.\\
The value of each bit can be:\\
0: Not set\\
1: The corresponding bit in \hyperref[regdesc:LPGPIOSTATUSREG]{LP\_GPIO\_STATUS\_INT} will be set\\
Bit[0]$\sim$bit[6] are corresponding to GPIO0$\sim$GPIO6.\\
Recommended operation: use this register to set \hyperref[regdesc:LPGPIOSTATUSREG]{LP\_GPIO\_STATUS\_INT}. \\
(WT)
\end{reglist}\end{regdesc}
\end{register}


\begin{register}{H}{LP\_GPIO\_STATUS\_W1TC\_REG}{0x0028}\label{regdesc:LPGPIOSTATUSW1TCREG}
\regfield{(reserved)}{25}{7}{0000000000000000000000000}%
\regfield{LP\_GPIO\_STATUS\_W1TC}{7}{0}{{0x0}}%
\reglabel{Reset}\regnewline%
\begin{regdesc}\begin{reglist}
\label{fielddesc:LPGPIOSTATUSW1TC}\item [LP\_GPIO\_STATUS\_W1TC] Configures whether or not to clear the interrupt status register \hyperref[regdesc:LPGPIOSTATUSREG]{LP\_GPIO\_STATUS\_INT} of GPIO0$\sim$GPIO6. \\
The value of each bit can be:\\
0: Not clear\\
1: The corresponding bit in \hyperref[regdesc:LPGPIOSTATUSREG]{LP\_GPIO\_STATUS\_INT} will be cleared\\
Bit[0]$\sim$bit[6] are corresponding to GPIO0$\sim$GPIO6.\\
Recommended operation: use this register to clear \hyperref[regdesc:LPGPIOSTATUSREG]{LP\_GPIO\_STATUS\_INT}.\\
(WT)
\end{reglist}\end{regdesc}
\end{register}


\begin{register}{H}{LP\_GPIO\_STATUS\_NEXT\_REG}{0x002C}\label{regdesc:LPGPIOSTATUSNEXTREG}
\regfield{(reserved)}{25}{7}{0000000000000000000000000}%
\regfield{LP\_GPIO\_STATUS\_INTERRUPT\_NEXT}{7}{0}{{0x0}}%
\reglabel{Reset}\regnewline%
\begin{regdesc}\begin{reglist}
\label{fielddesc:LPGPIOSTATUSINTERRUPTNEXT}\item [LP\_GPIO\_STATUS\_INTERRUPT\_NEXT] Represents the interrupt source status of GPIO0$\sim$GPIO6.\\
Bit[0]$\sim$bit[6] are corresponding to GPIO0$\sim$GPIO6.\\
Each bit represents:\\
0: Interrupt source status is invalid.\\
1: Interrupt source status is valid.\\ 
The interrupt here can be rising-edge triggered, falling-edge triggered, any edge triggered, or level triggered.\\ (RO)
\end{reglist}\end{regdesc}
\end{register}


\begin{register}{H}{LP\_GPIO\_PIN\regindex{n}\_REG (\regindex{n}: 0-6)}{0x0030+0x4*\regindex{n}}\label{regdesc:LPGPIOPINNREG}
\regfield{(reserved)}{21}{11}{000000000000000000000}%
\regfield{LP\_GPIO\_PIN\regindex{n}\_WAKEUP\_ENABLE}{1}{10}{{0}}%
\regfield{LP\_GPIO\_PIN\regindex{n}\_INT\_TYPE}{3}{7}{{0x0}}%
\regfield{(reserved)}{1}{6}{0}%
\regfield{LP\_GPIO\_PIN\regindex{n}\_EDGE\_WAKEUP\_CLR}{1}{5}{{0}}%
\regfield{LP\_GPIO\_PIN\regindex{n}\_SYNC1\_BYPASS}{2}{3}{{0x0}}%
\regfield{LP\_GPIO\_PIN\regindex{n}\_PAD\_DRIVER}{1}{2}{{0}}%
\regfield{LP\_GPIO\_PIN\regindex{n}\_SYNC2\_BYPASS}{2}{0}{{0x0}}%
\reglabel{Reset}\regnewline%
\begin{regdesc}\begin{reglist}
\label{fielddesc:LPGPIOPINNSYNC2BYPASS}\item [LP\_GPIO\_PIN\regindex{n}\_SYNC2\_BYPASS] Configures whether or not to synchronize GPIO input data on either edge of LP IO MUX operating clock for the second-level synchronization.\\
0: Not synchronize\\
1: Synchronize on falling edge\\
2: Synchronize on rising edge\\
3: Synchronize on rising edge\\ (R/W)
\label{fielddesc:LPGPIOPINNPADDRIVER}\item [LP\_GPIO\_PIN\regindex{n}\_PAD\_DRIVER] Configures to select the pin dirve mode of GPIO\regindex{n}.\\
0: Normal output\\
1: Open drain output\\ (R/W)
\label{fielddesc:LPGPIOPINNSYNC1BYPASS}\item [LP\_GPIO\_PIN\regindex{n}\_SYNC1\_BYPASS] Configures whether or not to synchronize GPIO input data on either edge of LP IO MUX operating clock for the first-level synchronization.\\
0: Not synchronize\\
1: Synchronize on falling edge\\
2: Synchronize on rising edge\\
3: Synchronize on rising edge\\ (R/W)
\label{fielddesc:LPGPIOPINNEDGEWAKEUPCLR}\item [LP\_GPIO\_PIN\regindex{n}\_EDGE\_WAKEUP\_CLR] Configures whether or not to clear the edge wake-up status of GPIO0$\sim$GPIO6.\\
0: No effect\\
1: Clear\\
(WT)
\label{fielddesc:LPGPIOPINNINTTYPE}\item [LP\_GPIO\_PIN\regindex{n}\_INT\_TYPE] Configures GPIO\regindex{n} interrupt type.\\
    0: GPIO interrupt disabled \\
    1: Rising edge trigger \\
    2: Falling edge trigger \\
    3: Any edge trigger \\
    4: Low level trigger \\
    5: High level trigger \\
 (R/W)
\item [Continued on the next page...]
\end{reglist}\end{regdesc}
\end{register}

\addtocounter{Regfloat}{-1}
\begin{register}{H}{LP\_GPIO\_PIN\regindex{n}\_REG (\regindex{n}: 0-6)}{0x0030+0x4*\regindex{n}}
\begin{regdesc}\begin{reglist}
\item [Continued from the previous page...]
\label{fielddesc:LPGPIOPINNWAKEUPENABLE}\item [LP\_GPIO\_PIN\regindex{n}\_WAKEUP\_ENABLE] Configures whether or not to enable GPIO\regindex{n} wake-up function.\\
0: Not enable\\
1: Enable\\
This function is disabled when PD\_LP\_PERI is powered off. For more information, see Chapter \nameref{mod:lowpowm}.\\ (R/W)
\end{reglist}\end{regdesc}
\end{register}


\begin{register}{H}{LP\_GPIO\_FUNC\regindex{n}\_OUT\_SEL\_CFG\_REG (\regindex{n}: 0-6)}{0x02B0+0x4*\regindex{n}}\label{regdesc:LPGPIOFUNCNOUTSELCFGREG}
\regfield{(reserved)}{29}{3}{00000000000000000000000000000}%
\regfield{LP\_GPIO\_FUNC\regindex{n}\_OE\_INV\_SEL}{1}{2}{{0}}%
\regfield{(reserved)}{1}{1}{0}%
\regfield{LP\_GPIO\_FUNC\regindex{n}\_OUT\_INV\_SEL}{1}{0}{{0}}%
\reglabel{Reset}\regnewline%
\begin{regdesc}\begin{reglist}
\label{fielddesc:LPGPIOFUNCNOUTINVSEL}\item [LP\_GPIO\_FUNC\regindex{n}\_OUT\_INV\_SEL] Configures whether or not to invert the output value.\\
0: Not invert\\
1: Invert\\ (R/W)
\label{fielddesc:LPGPIOFUNCNOEINVSEL}\item [LP\_GPIO\_FUNC\regindex{n}\_OE\_INV\_SEL] Configures whether or not to invert the output enable signal.\\
0: Not invert\\
1: Invert\\ (R/W)
\end{reglist}\end{regdesc}
\end{register}


\begin{register}{H}{LP\_GPIO\_CLOCK\_GATE\_REG}{0x03F8}\label{regdesc:LPGPIOCLOCKGATEREG}
\regfield{(reserved)}{31}{1}{0000000000000000000000000000000}%
\regfield{LP\_GPIO\_CLK\_EN}{1}{0}{{1}}%
\reglabel{Reset}\regnewline%
\begin{regdesc}\begin{reglist}
\label{fielddesc:LPGPIOCLKEN}\item [LP\_GPIO\_CLK\_EN] Configure to enable 
 the GPIO clock gate or not.\\
0: Not enable\\
1: Enable \\ 
(R/W)
\end{reglist}\end{regdesc}
\end{register}


\begin{register}{H}{LP\_GPIO\_DATE\_REG}{0x03FC}\label{regdesc:LPGPIODATEREG}
\regfield{(reserved)}{4}{28}{0000}%
\regfield{LP\_GPIO\_DATE}{28}{0}{{0x2310240}}%
\reglabel{Reset}\regnewline%
\begin{regdesc}\begin{reglist}
\label{fielddesc:LPGPIODATE}\item [LP\_GPIO\_DATE] Version control register.\\ (R/W)
\end{reglist}\end{regdesc}
\end{register}


